% Resumo (NBR 6028: um parágrafo, 150-500 palavras, sem fórmulas).
% Reescrito em 24/11/2026 (revisão pedagógica); todos os números conferidos
% nos capítulos e artefatos de origem antes da inclusão.
Por uma década, a escolha de modelos para dados tabulares apoiou-se na
máxima de que o \emph{gradient boosting} sempre vence. Os
\emph{foundation models} tabulares, redes pré-treinadas que predizem sem
treinamento na tarefa, puseram essa máxima em dúvida, mas a evidência
disponível vem sobretudo dos próprios proponentes, sob protocolos
distintos. Esta dissertação tem por objetivo produzir evidência
comparativa sob protocolo único e convertê-la em um procedimento de
escolha, em três frentes. Primeiro, um benchmark multicritério avalia 14
modelos de três famílias (\emph{gradient boosting}, aprendizado profundo
e \emph{foundation models}) em 18 conjuntos públicos do OpenML, com
\emph{hold-out} que nenhuma etapa de seleção observa, inferência
frequentista e bayesiana e medição dos custos de treino, de ajuste e de
latência. Segundo, um arcabouço de decisão é avaliado: o componente aprendido,
deixando um
conjunto de fora por vez (\emph{leave-one-dataset-out}); as políticas
fixas, descritivamente, nos mesmos conjuntos. Terceiro,
estuda-se a destilação das curvas de quantis do professor TabPFN em
alunos de \emph{gradient boosting}, em 20 conjuntos de regressão, sob
hipóteses registradas. Entre os 11 modelos anteriores a 2026, o topo da
acurácia é um empate: de 3 a 8 dos 10 modelos não líderes são
praticamente equivalentes ao líder, conforme o critério bayesiano
adotado, na margem de um ponto percentual de AUC. Na geração 2026, o
TabFM, sem ajuste algum, tem o melhor rank médio nas três tarefas (na
multiclasse, sobre os dois conjuntos em que opera) e é o
melhor modelo em 13 dos 18 conjuntos; a liderança é consistente entre
conjuntos, mas os testes, de baixo poder com tão poucos conjuntos por
tarefa, não a confirmam como separação estatística, e ela custa de
quatro a cinco ordens de magnitude a mais em latência de inferência. As
redes de Kolmogorov--Arnold ficam no último terço nas três tarefas. Na
decisão, um roteador aprendido por meta-características acerta menos que
a regra trivial (0,11 contra 0,44); a política fixa ``\emph{foundation
model} primeiro'' tem \emph{regret} mediano zero na métrica primária,
mas deixa de ser a recomendada quando 10 a 20 por cento do peso da
decisão passa da acurácia para a latência, na normalização por postos
adotada. Na destilação, o aluno
treinado nas médias do professor supera o controle em apenas 4 de 20
conjuntos; o aluno treinado nas curvas de quantis retém, na mediana, de
12 a 19 por cento da vantagem distribucional do professor, conforme o
ajuste dos alunos, com retenção positiva em 16 de 20 conjuntos (teste de
sinal, p de aproximadamente 0,019 com a correção do desenho em duas
etapas), e supera \emph{baselines} distribucionais clássicos em 17 de
20. Conclui-se que, no regime estudado (conjuntos de até 100 mil
amostras, versões fixadas dos modelos, uma única GPU), a escolha deve
partir das restrições da aplicação (latência, limites arquiteturais,
interpretabilidade), e não de um vencedor universal; e que a destilação
recupera uma parte modesta, porém consistente, da incerteza preditiva do
\emph{foundation model}, sob pré-condições que se verificam em minutos.

\vspace{1cm}
\noindent\textbf{Palavras-chave:} dados tabulares; foundation models;
benchmark; destilação de conhecimento; regressão distribucional.
